Die primäre Zielgruppe dieser Arbeit sind Entwickler, die an der Konzeption und Implementierung KI-gestützer Textverarbeitungssysteme beteiligt sind, insbesondere im Bereich der barrierefreien Kommunikation. Darüber hinaus richtet sich die Arbeit an Sprachwissenschaftler und Forschende im Bereich Verständlichkeitsforschung sowie an Medienhäuser und öffentlich-rechtliche Rundfunkanstalten, die automatisierte Verfahren zur Erstellung von Leichter Sprache einsetzen oder planen.

Sekundär adressiert die Arbeit auch Studenten und Frachkräfte, die sich mit Large Language Models, Textbewertung und iterativen Optimierungsverfahren beschäftigen, da die vorgestellten Konzepte und Methoden auch in anderen Kontexten der natürlichen Sprachverarbeitung Anwendung finden können.